\section{Introduction}
\label{sec:introduction}

\subsection{From counting to closure}
A familiar way to ``count a stone'' is to apply a tiny, discrete action many times: tick marks, steps along a line, repeated updates of a running total.
A familiar way to \emph{do calculus} is to replace that large discrete composition by a compressed law: a differential equation for the limit of many tiny updates, or an integral for the limit of many tiny accumulations.
This paper takes that replacement seriously as a \emph{closure problem}.
The central question is not ``how do we define derivatives?'' but:
\emph{which compressions of discrete protocols remain coherent as we refine the stage parameter toward a continuum limit?}

In the \SBT dictionary (reviewed in Section~2), refinement is \Pfour\ staging; the defect budget carried by a compression is \Psix\ accounting; and the limit object is \Pfive\ packaging.
On this view, the basic moves of analysis are forced by feasibility: only those operators and equivalence relations whose defects stabilize under refinement become admissible ``laws'' at the packaged layer.
The same lens applies beyond calculus, whenever one attempts to compress a complicated discrete structure into a continuous or spectral object.

This manuscript is a substantially revised version of the Zenodo preprint (DOI: \href{https://doi.org/10.5281/zenodo.18402004}{10.5281/zenodo.18402004}); see Related Identifiers for version linkage.

\subsection{Paper contract: what this paper is (and is not)}
This document is a concrete instantiation of the general \SBT framework of \cite{SixBirdsFramework} in the domain of mathematics, and it treats the core toolkit as an \emph{emergence calculus} for stable closure under refinement.

\paragraph{This paper \emph{is}:}
\begin{itemize}
  \item a worked translation of the \SBT primitives into standard mathematical practice: limits, completions, renormalizations, and compatibility constraints;
  \item a validation-first report in which key claims are supported either by machine-checked lemmas (Lean) or by falsification-oriented numerical diagnostics (Python);
  \item a proposal for how ``higher'' mathematical layers can be treated as \emph{closures} of discrete protocol algebras under refinement, rather than as isolated definitions.
\end{itemize}

\paragraph{This paper \emph{is not}:}
\begin{itemize}
  \item a replacement for standard foundations or a proposal of new axioms;
  \item a proof of the Riemann Hypothesis or a claim that the diagnostics in Exhibits~B-1/B-2 settle classical conjectures;
  \item an attempt to reduce all of mathematics to a single mechanism.
\end{itemize}

The intended reading is pragmatic: \SBT is used here as a \emph{dictionary} for organizing constructions and for designing tests that can fail.
When the diagnostics fail (as they sometimes do), we treat the failure as information about missing feasibility conditions rather than as a defect to be hidden.

\subsection{Contributions and artifacts}
The paper contributes a small set of concrete, reusable ``exhibits'':

\begin{itemize}
  \item \textbf{Exhibit A-1 (calculus closure).} We isolate the exact discrete Leibniz identity for finite differences (with its explicit remainder term) and use it as an accounting ledger: demanding that the remainder vanish in a packaged limit forces the derivation (product rule) structure. We also record an algebraic uniqueness anchor for derivations on polynomial rings in Lean/mathlib \cite{moura2021lean4,mathlib2020lean}.
  \item \textbf{Exhibit A-2 (protocol holonomy).} We quantify route mismatch between two natural refinement protocols under a small coordinate change, treating the mismatch as a \Pthree\ holonomy-style defect that decays with refinement.
  \item \textbf{Exhibit B-1 (prime closure diagnostics).} We compare staged additive and multiplicative closures (Dirichlet truncations versus Euler-product truncations) under prototype packaging maps, and we separate a convergence control regime from the critical strip where naive staging is unstable.
  \item \textbf{Exhibit B-2 (passivity toy).} We show, in a self-dual toy family with a positivity constraint, that tightening feasibility can confine zeros to a symmetry locus---a pattern that motivates (but does not prove) ``passivity ledger'' interpretations of zero confinement.
\end{itemize}

A key implementation choice is that reported numbers and tables are not hand-copied: the TeX paper inputs generated tables and macros produced from snapshot-visible ``last run'' pointers in \texttt{notes/}.
This makes the draft internally consistent with the repository state.

\paragraph{Code availability.}
The math-instantiation repository is available at:
\begin{itemize}[noitemsep,topsep=2pt]
\item \url{https://github.com/ioannist/six-birds-mathematics}
\end{itemize}

\subsection{Roadmap}
Section~2 recalls the \SBT dictionary in the minimal form needed for this instantiation.
Section~3 formulates the discrete-to-packaged ``closure engine'' for mathematics.
Sections~4--7 present the four exhibits: A-1 (calculus closure), A-2 (protocol holonomy), B-1 (prime closure diagnostics), and B-2 (passivity toy).
Section~8 documents the artifact pipeline and reproduction commands, and Section~9 discusses limitations and next steps.
